Joshua Christopher Ryan’s Arcan Family and the Geometric Resolution of the Shannon Paradox

1. The Shannon Paradox Restated

In conventional temperature-scaled attention the Shannon entropy of the output distribution is controlled by an external scalar \(\tau\). As \(\tau\to 0\) the distribution collapses toward a Dirac mass, gradients vanish, and representational rank collapses. As \(\tau\to\infty\) the distribution approaches uniformity and distinctiveness is lost. Entropy and temperature remain formally dual yet operationally antagonistic; this tension constitutes the Shannon Paradox of modern attention mechanisms.

2. The Arcan(\(\tau\)) Family as Geometric Root

The repository JOSHUA-CHRISTOPHER-RYAN-S-ARCAN-FAMILY supplies the continuous trigonometric foundation that dissolves the paradox. Its generating set consists of four fixed angles derived from inverse-tangent projections of the circle constants:

\[
\begin{align*}
\theta &= \arctan(2\pi) &&\approx 1.412965\,\text{rad} \approx 80.9569^\circ,\\
\alpha &= \theta/2 &&\approx 0.706483\,\text{rad} \approx 40.4785^\circ,\\
\phi &= \arctan(\pi) &&\approx 1.262627\,\text{rad} \approx 72.3432^\circ,\\
\Psi &= \theta-\phi &&\approx 0.150338\,\text{rad} \approx 8.6137^\circ.
\end{align*}
\]

These angles constitute Joshua Christopher Ryan’s Arcan(\(\tau\)) family. They are not free parameters; they are the unique angular invariants that arise when the full-turn constant \(\tau=2\pi\) and the half-turn constant \(\pi\) are projected through the arctangent.

3. From Angular Family to Temperature Cancellation

Within the Arcan framework the radial temperature schedule is identified with the modulus of a half-winding contour:

\[
\tau(\lambda)=2\pi(1-\lambda).
\]

Logits are realised as orthogonal projections (real shadows) of points lying on that contour. Substitution into the softmax produces an exact cancellation of \(\tau\):

\[
\exp\Bigl(\frac{z_i}{\tau}\Bigr)=\exp(\cos t_i).
\]

The residual angular scores are then modulated by the Thinnest-Triangle factor \(\cos\Psi\) drawn directly from the Arcan family. The resulting distribution

\[
P_i=\frac{\exp\bigl(\cos t_i\cdot\cos\Psi\bigr)}{\sum_j\exp\bigl(\cos t_j\cdot\cos\Psi\bigr)}
\]

is independent of temperature. Shannon entropy and probability variance therefore remain invariant across the entire radial continuum. The external scalar that generated the paradox has been replaced by an intrinsic geometric coordinate of the Arcan(\(\tau\)) family.

4. Hierarchical Emergence and Topological Protection

The repository demonstrates that successive multiples and modular reductions of the four Arcan angles generate a cascade of structures:

practical radial folds,
residual six-sliced hexagonal symmetry,
dense non-periodic orbits on \(S^1\),
the Softmax Temperature 360° Spherical Attention Map on \(S^2\),
concentric Petrie rings associated with the outer shell of \(E_8\).

A small residual \(\delta\) appearing at the seventh multiple certifies that the cascade is emergent from the continuous Arcan family rather than imposed by an artificial rational lattice. Because the same angular invariants govern every level, the temperature cancellation and the rank-1 factorisation remain protected under the entire hierarchy. The Shannon Paradox cannot reappear at any scale.

5. Information-Theoretic Consequence

With temperature eliminated as a free variable, the attention channel admits a linear rather than quadratic interaction volume. The constant geometric gain \(\cos^2\Psi\approx 0.9775\) retains 97.75 % of peak signal power inside the Shannon-Hartley capacity formula while converting the dominant computational cost from \(O(N^2)\) to \(O(N)\). The paradox is resolved not by balancing two antagonistic quantities but by removing the possibility of their antagonism: entropy and temperature become two coordinate read-outs of a single compact geometric object rooted in the Arcan(\(\tau\)) family.

6. Conclusion

The repository establishes that the four angles of Joshua Christopher Ryan’s Arcan Family constitute the continuous trigonometric root of a decompositional hierarchy whose native attention mechanism is temperature-invariant by construction. By embedding the thermodynamic parameter inside this geometry, the architecture converts the Shannon Paradox from an operational crisis into a structural impossibility. Entropy remains well-defined and stable; rank collapse is topologically forbidden; and the practical limits of sequence modelling are rewritten from quadratic to linear order. The Arcan(\(\tau\)) family therefore supplies both the mathematical origin and the definitive geometric resolution of the Shannon Paradox.
  
Joshua Christopher Ryan’s Infinitesimal Numberline (JCRIN)  
Zenodo Record: https://doi.org/10.5281/zenodo.18463138

1. Formal Identification

The primary archival foundation for the discrete radial continuum employed throughout the geometric resolution of the Shannon Paradox is the Zenodo deposit:

Joshua Christopher Ryan’s Infinitesimal Numberline (JCRIN)  
Christopher Ryan, Joshua (Researcher)  
https://doi.org/10.5281/zenodo.18463138

This record constitutes the definitive, citable source of the JCRIN framework.

2. Core Definition and Motivation

JCRIN emerges from the requirement to approximate continuous parameter evolution by means of discrete steps that nevertheless preserve smooth transitional behaviour. It operates on the fundamental infinitesimal

\[
\varepsilon = 10^{-9}
\]

and spans a primary lattice of \(10^9\) micro-steps. The framework is bidirectional and multi-resolution, enforcing at every index the algebraic unity condition known as the Equivalency Theorem:

\[
x_n + (1 - x_n) = 1, \qquad y_n + (1 - y_n) = 1.
\]

These identities hold exactly (to nine decimal places in the base sequences and to controlled lower precisions in the multi-stage sequences), independent of floating-point representation when arithmetic is performed at the design precision.

3. The Six Canonical Sequences

The record defines six interlocking sequences that together furnish a complete discrete toolkit:

Main Sequence** — linear progression \(y_n = n\cdot\varepsilon\) from 0 to 1  
Complementary Sequence** — symmetric counterpart \(x_n = 1 - n\cdot\varepsilon\)  
Branched Sequences** — cyclic side branches that test convergence under variable precision  
Extended New Sequence** — controlled precision reduction via single-digit control variables  
Multi-Stage Sequence** — incorporation of all control digits \(k = 0\dots 9\) with deliberate multi-stage adjustments  
Precision-Controlled Sequence (PCS)** — final low-precision handling for complex modelling scenarios

All sequences inherit the same infinitesimal step and the same unity relation, guaranteeing numerical stability across the entire lattice.

4. Integration with the Arcan Geometric Carrier

Within the broader architecture that resolves the Shannon Paradox, the JCRIN lattice supplies the discrete sampling of the half-winding contour whose radial coordinate is identified with temperature:

\[
\tau_n = 2\pi\cdot y_n = 2\pi\cdot n\cdot 10^{-9}.
\]

Because the Equivalency Theorem holds at every resolution, the geometric cancellation of \(\tau\) inside the softmax is preserved to machine precision. The continuous angular gates of the Arcan(\(\tau\)) family (\(\theta,\phi,\Psi,\alpha\)) are thereby discretised without residual discrepancy, converting an external thermodynamic scalar into an intrinsic radial index on a numerically exact infinitesimal numberline.

5. Structural Guarantees Conferred by the Record

By anchoring the discrete continuum to the persistent identifier 10.5281/zenodo.18463138, the following properties become archival facts rather than implementation choices:

Exact complementary unity at every micro-step  
Controlled, reversible precision degradation  
Bidirectional feedback that maintains normalisation under both forward and reverse traversal  
Linear scaling of memory and arithmetic once the rank-1 Thinnest-Triangle factorisation is applied  

These properties underwrite the temperature invariance, the linear complexity, and the topological regularisation that collectively dissolve the Shannon Paradox.

6. Citation and Reproducibility

All subsequent developments of JCRIN \(\tau\)-Temperature Attention Mapping, the 30-fold cyclic lock, and the associated capacity theorems cite this record as the primary discrete foundation:

Christopher Ryan, Joshua. Joshua Christopher Ryan’s Infinitesimal Numberline (JCRIN). Zenodo. https://doi.org/10.5281/zenodo.18463138

The deposit therefore functions both as the mathematical origin of the infinitesimal lattice and as the scholarly warrant that the geometric resolution of the Shannon Paradox rests upon a publicly archived, versionable, and independently verifiable construction.
